\documentclass[11pt]{amsart}
\ifdefined\XeTeXrevision\else\ifdefined\pdfoutput\pdfoutput=1\fi\fi

\usepackage[margin=1.2in]{geometry}
\usepackage{amsmath,amssymb,mathtools}
\usepackage{booktabs}
\usepackage{array}
\usepackage{xcolor}
\definecolor{linkblue}{RGB}{25,60,120}
\usepackage[colorlinks=true,linkcolor=linkblue,citecolor=linkblue,urlcolor=linkblue]{hyperref}
\hypersetup{
  pdftitle={Peaceable queens on the 15x15 torus: t(15)=20},
  pdfauthor={Jan Philipp Harries},
  pdfsubject={Peaceable queens on toroidal chessboards},
  pdfkeywords={peaceable queens, toroidal chessboard, exhaustive search, computer-assisted proof, OEIS A279405}
}

\newtheorem{theorem}{Theorem}
\newtheorem{lemma}[theorem]{Lemma}
\newtheorem{proposition}[theorem]{Proposition}
\newtheorem{corollary}[theorem]{Corollary}
\theoremstyle{remark}
\newtheorem{remark}[theorem]{Remark}

\newcommand{\Z}{\mathbb{Z}}
\newcommand{\Zn}{\Z/15\Z}
\newcommand{\Bl}{\mathrm{B}}
\newcommand{\Wh}{\mathrm{W}}

\title[Peaceable queens on the $15\times 15$ torus]{Peaceable queens on the $15\times 15$ torus:\\ the exact value $t(15)=20$}

\author{Jan Philipp Harries}
\thanks{ellamind GmbH, Bremen, Germany. \textit{E-mail address:} \href{mailto:jan@ellamind.com}{jan@ellamind.com}}

\date{July 25, 2026}

\subjclass[2020]{Primary 05B99; Secondary 05C69, 68R05, 05C30}
\keywords{peaceable queens, toroidal chessboard, exhaustive search, computer-assisted proof, OEIS A279405}

\begin{document}

\begin{abstract}
Let $t(n)$ denote the largest $m$ such that $m$ black and $m$ white queens can be placed on the $n\times n$ toroidal chessboard with no black queen attacking a white queen. We prove that $t(15)=20$, determining the first previously unknown term of this sequence (OEIS A279405), whose exact values were known only through $n=14$. The lower bound is an explicit $20+20$ placement. The upper bound is a computer-assisted exhaustive proof: a queen placement is replaced by a colouring of the $60$ toroidal lines, elementary counting identities restrict the four line-set cardinalities to $247$ profiles up to a symmetry group of order $14{,}400$, and each profile is decided exactly by a finite search whose every pruning step is a necessary condition. The search examined $6{,}074{,}753{,}568$ cases and found no $21+21$ placement. Two independently written exact enumerators, using different symmetry quotients and different completion algorithms, agree on all $247$ profiles.
\end{abstract}

\maketitle

\section{Introduction}\label{sec:intro}

A \emph{peaceable} pair of queen armies is a placement of $m$ black and $m$ white queens on a chessboard such that no black queen attacks a white queen; queens of the same colour may attack each other freely. The problem of maximizing $m$ goes back to Ainley \cite{Ainley77} and was popularized through the OEIS sequence A250000 for the ordinary $n\times n$ board. On the \emph{toroidal} board the four lines through a cell wrap around, which makes the problem homogeneous: the symmetry group acts transitively on cells, and the combinatorics becomes arithmetic in $\Zn$. The toroidal maxima form OEIS sequence A279405 \cite{OEIS279405}.

\begin{table}[h]
\centering
\begin{tabular}{@{}l*{8}{r}@{}}
\toprule
$n$      & 8 & 9 & 10 & 11 & 12 & 13 & 14 & \textbf{15}\\
$t(n)$   & 8 & 7 & 12 & 10 & 18 & 16 & 25 & \textbf{20}\\
$n^2/8$  & 8.0 & 10.1 & 12.5 & 15.1 & 18.0 & 21.1 & 24.5 & 28.1\\
\bottomrule
\end{tabular}
\medskip
\caption{The upper end of A279405. The values through $n=14$ are those recorded in \cite{OEIS279405}; $a(13)=16$ and $a(14)=25$ were contributed by Jaideep Padhi in June 2026. The value at $n=15$ is Theorem~\ref{thm:main}.}\label{tab:oeis}
\end{table}

Table~\ref{tab:oeis} shows the characteristic parity split of the toroidal problem: for even $n\le 14$ the known values sit close to $n^2/8$, whereas for odd $n$ they fall well short of it, hovering nearer $n^2/11$. Clinch, Drescher, Huynh and Saffidine \cite[Theorem~1.4]{CDHS24} proved the bound $t(n)\le n^2/8$ for odd $n$, which for $n=15$ gives $t(15)\le 28$ by integrality, and exhibited an explicit $20+20$ configuration for $n=15$ \cite[Appendix~A]{CDHS24}. The value $t(15)$ was therefore known to lie in $[20,28]$, and A279405 listed no term for $n=15$.

We determine it.

\begin{theorem}\label{thm:main}
On the $15\times 15$ toroidal chessboard, $20$ black and $20$ white queens can be placed with no black--white attack, and $21$ black and $21$ white queens cannot. That is,
\[
t(15)=20 .
\]
\end{theorem}

The lower bound is a placement, reproduced in Section~\ref{sec:lower}. The upper bound is the substance of this note and is computer-assisted. Its structure is as follows.

Section~\ref{sec:reduction} eliminates the queens. A $21+21$ placement exists if and only if the $60$ lines of the board can be coloured black and white so that the set of cells all four of whose lines are black, and the set of cells all four of whose lines are white, both have at least $21$ elements. The black line sets $R,C,D,A\subseteq\Zn$ (rows, columns, diagonals, antidiagonals) then determine everything, and swapping colours if necessary one may assume $S:=|R|+|C|+|D|+|A|\le 30$. We stress that $S\le 30$ is genuinely an inequality: it cannot be normalized to $S=30$, because adding a line to the black side removes it from the white side and may destroy the white count. Indeed the placement of Section~\ref{sec:lower} has $S=28$.

Section~\ref{sec:profiles} bounds the possible \emph{size profiles} $(|R|,|C|,|D|,|A|)$. Two lines from different families meet in exactly one cell --- for the diagonal--antidiagonal pair this uses the invertibility of $2$ modulo $15$ --- which yields the pair conditions $s_is_j\ge 21$ and $(15-s_i)(15-s_j)\ge 21$. Expanding three complementary indicator functions yields an exact identity whose consequence is $225-15(x+y+z)+xy+xz+yz\ge 42$ for each triple of family sizes. The affine symmetries of the board that permute the four line families form a group of order $14{,}400$ inducing a dihedral action of order $8$ on the families; together with colour complementation in the stratum $S=30$ this leaves exactly $247$ canonical profiles.

Section~\ref{sec:completion} decides a single profile exactly. One fixes a pair of families up to the symmetry quotient, enumerates all subsets of one remaining family, and decides the last family by an exact $15$-item, fixed-cardinality, two-constraint subset problem. Every pruning rule used is a necessary condition, so the procedure is exhaustive and not a heuristic search.

Section~\ref{sec:result} reports the outcome: all $247$ profiles are infeasible, over $6{,}074{,}753{,}568$ enumerated cases. Section~\ref{sec:verification} describes the verification, which deserves the most scrutiny and accordingly receives the most space: two independently written exact enumerators with different quotient strategies and different completion algorithms, three runs of the proof solver on two toolchains, a Burnside-lemma audit of every orbit count, sanitizer builds, and brute-force gates on the completion kernel. Sections~\ref{sec:repro} and~\ref{sec:ai} give reproduction instructions and disclose the use of AI systems in producing this note.

\begin{remark}
The proof is exhaustive rather than explanatory. It gives no human-readable reason why the odd torus is so much poorer than the even torus, and it does not sharpen the bound of \cite{CDHS24} for general odd $n$. A structural proof of $t(15)\le 20$ --- or, better, of a general odd-$n$ bound near $n^2/11$ --- remains open and would supersede this computation.
\end{remark}

\section{Notation}\label{sec:notation}

Throughout, $n=15$ and all coordinates lie in $\Zn$, represented by $\{0,1,\dots,14\}$. The board is $G=(\Zn)^2$, with $225$ cells. Through a cell $(r,c)$ pass four \emph{lines}:
\[
\text{row } r,\qquad \text{column } c,\qquad \text{diagonal } d=r-c,\qquad \text{antidiagonal } a=r+c ,
\]
each family indexed by $\Zn$, so the board carries $60$ lines. Two queens \emph{attack} each other exactly when they share a line. Since $2$ is invertible modulo $15$, with $2^{-1}=8$, the map $(r,c)\mapsto(d,a)=(r-c,r+c)$ is a bijection of $G$, with inverse
\begin{equation}\label{eq:inverse}
r=8(d+a),\qquad c=8(a-d).
\end{equation}
In particular a diagonal and an antidiagonal meet in exactly one cell, just as a row and a column do; this is the one place where the arithmetic of $\Zn$ (rather than of an arbitrary $\Z/n\Z$) is used, and it fails for even $n$.

We write $t(n)$ for the largest $m$ such that $m$ black and $m$ white queens can be placed on the $n\times n$ torus with no black--white attack. All $2m$ queens occupy distinct cells; queens of one colour may share lines.

\section{The lower bound}\label{sec:lower}

\begin{proposition}\label{prop:lower}
$t(15)\ge 20$.
\end{proposition}

\begin{proof}
Take the black line sets
\begin{equation}\label{eq:witness}
\begin{aligned}
R&=\{5,7,8,11,12,14\}, & D&=\{0,1,4,8,12,13,14\},\\
C&=\{0,3,4,7,9,10,12,13\}, & A&=\{1,2,3,9,10,12,14\},
\end{aligned}
\end{equation}
of sizes $(6,8,7,7)$, and colour every remaining line white, giving white line sets of sizes $(9,7,8,8)$. Call a cell \emph{black-homogeneous} if $r\in R$, $c\in C$, $r-c\in D$ and $r+c\in A$, and \emph{white-homogeneous} if $r\notin R$, $c\notin C$, $r-c\notin D$ and $r+c\notin A$. Direct enumeration of the $225$ cells gives exactly $20$ black-homogeneous cells,
\[
\begin{array}{llll}
(5,4)&(5,7)&(5,12)&\\
(7,3)&(7,7)&(7,9)&(7,10)\\
(8,4)&(8,9)&(8,10)&\\
(11,3)&(11,7)&(11,13)&\\
(12,0)&(12,4)&(12,12)&(12,13)\\
(14,0)&(14,10)&(14,13)&
\end{array}
\]
and exactly $22$ white-homogeneous cells, of which we use the following $20$:
\[
\begin{array}{llll}
(0,5)&(0,6)&(0,8)&\\
(1,5)&(1,6)&(1,14)&\\
(2,6)&(2,11)&&\\
(3,1)&(3,8)&&\\
(4,1)&(4,2)&&\\
(6,1)&(6,14)&&\\
(9,2)&(9,6)&(9,14)&\\
(10,1)&(10,5)&(13,2)&
\end{array}
\]
(the two unused ones are $(13,6)$ and $(13,8)$). Place black queens on the first list and white queens on the second. Every black queen lies on four black lines and every white queen on four white lines; a black--white attack would require a line that is both black and white, and there is none. Hence $20$ black and $20$ white queens coexist peaceably.
\end{proof}

Figure~\ref{fig:board} displays the configuration. This lower bound is not new: an explicit $20+20$ configuration for $n=15$ is given in the appendix of \cite{CDHS24}. We reproduce a placement here in line-set form so that Theorem~\ref{thm:main} is self-contained and checkable by hand.

\begin{figure}[t]
\centering
\setlength{\tabcolsep}{2.4pt}
\renewcommand{\arraystretch}{1.05}
{\small
\begin{tabular}{r|*{15}{c}}
$r\backslash c$ & 0 & 1 & 2 & 3 & 4 & 5 & 6 & 7 & 8 & 9 & 10 & 11 & 12 & 13 & 14 \\
\midrule
0 & $\cdot$ & $\cdot$ & $\cdot$ & $\cdot$ & $\cdot$ & $\square$ & $\square$ & $\cdot$ & $\square$ & $\cdot$ & $\cdot$ & $\cdot$ & $\cdot$ & $\cdot$ & $\cdot$ \\
1 & $\cdot$ & $\cdot$ & $\cdot$ & $\cdot$ & $\cdot$ & $\square$ & $\square$ & $\cdot$ & $\cdot$ & $\cdot$ & $\cdot$ & $\cdot$ & $\cdot$ & $\cdot$ & $\square$ \\
2 & $\cdot$ & $\cdot$ & $\cdot$ & $\cdot$ & $\cdot$ & $\cdot$ & $\square$ & $\cdot$ & $\cdot$ & $\cdot$ & $\cdot$ & $\square$ & $\cdot$ & $\cdot$ & $\cdot$ \\
3 & $\cdot$ & $\square$ & $\cdot$ & $\cdot$ & $\cdot$ & $\cdot$ & $\cdot$ & $\cdot$ & $\square$ & $\cdot$ & $\cdot$ & $\cdot$ & $\cdot$ & $\cdot$ & $\cdot$ \\
4 & $\cdot$ & $\square$ & $\square$ & $\cdot$ & $\cdot$ & $\cdot$ & $\cdot$ & $\cdot$ & $\cdot$ & $\cdot$ & $\cdot$ & $\cdot$ & $\cdot$ & $\cdot$ & $\cdot$ \\
5 & $\cdot$ & $\cdot$ & $\cdot$ & $\cdot$ & $\blacksquare$ & $\cdot$ & $\cdot$ & $\blacksquare$ & $\cdot$ & $\cdot$ & $\cdot$ & $\cdot$ & $\blacksquare$ & $\cdot$ & $\cdot$ \\
6 & $\cdot$ & $\square$ & $\cdot$ & $\cdot$ & $\cdot$ & $\cdot$ & $\cdot$ & $\cdot$ & $\cdot$ & $\cdot$ & $\cdot$ & $\cdot$ & $\cdot$ & $\cdot$ & $\square$ \\
7 & $\cdot$ & $\cdot$ & $\cdot$ & $\blacksquare$ & $\cdot$ & $\cdot$ & $\cdot$ & $\blacksquare$ & $\cdot$ & $\blacksquare$ & $\blacksquare$ & $\cdot$ & $\cdot$ & $\cdot$ & $\cdot$ \\
8 & $\cdot$ & $\cdot$ & $\cdot$ & $\cdot$ & $\blacksquare$ & $\cdot$ & $\cdot$ & $\cdot$ & $\cdot$ & $\blacksquare$ & $\blacksquare$ & $\cdot$ & $\cdot$ & $\cdot$ & $\cdot$ \\
9 & $\cdot$ & $\cdot$ & $\square$ & $\cdot$ & $\cdot$ & $\cdot$ & $\square$ & $\cdot$ & $\cdot$ & $\cdot$ & $\cdot$ & $\cdot$ & $\cdot$ & $\cdot$ & $\square$ \\
10 & $\cdot$ & $\square$ & $\cdot$ & $\cdot$ & $\cdot$ & $\square$ & $\cdot$ & $\cdot$ & $\cdot$ & $\cdot$ & $\cdot$ & $\cdot$ & $\cdot$ & $\cdot$ & $\cdot$ \\
11 & $\cdot$ & $\cdot$ & $\cdot$ & $\blacksquare$ & $\cdot$ & $\cdot$ & $\cdot$ & $\blacksquare$ & $\cdot$ & $\cdot$ & $\cdot$ & $\cdot$ & $\cdot$ & $\blacksquare$ & $\cdot$ \\
12 & $\blacksquare$ & $\cdot$ & $\cdot$ & $\cdot$ & $\blacksquare$ & $\cdot$ & $\cdot$ & $\cdot$ & $\cdot$ & $\cdot$ & $\cdot$ & $\cdot$ & $\blacksquare$ & $\blacksquare$ & $\cdot$ \\
13 & $\cdot$ & $\cdot$ & $\square$ & $\cdot$ & $\cdot$ & $\cdot$ & $\lozenge$ & $\cdot$ & $\lozenge$ & $\cdot$ & $\cdot$ & $\cdot$ & $\cdot$ & $\cdot$ & $\cdot$ \\
14 & $\blacksquare$ & $\cdot$ & $\cdot$ & $\cdot$ & $\cdot$ & $\cdot$ & $\cdot$ & $\cdot$ & $\cdot$ & $\cdot$ & $\blacksquare$ & $\cdot$ & $\cdot$ & $\blacksquare$ & $\cdot$ \\
\end{tabular}}
\caption{The $20+20$ placement of Proposition~\ref{prop:lower}: $\blacksquare$ black queens, $\square$ white queens, $\lozenge$ the two white-homogeneous cells left empty. Rows, columns and both diagonal families wrap around.}\label{fig:board}
\end{figure}

\section{The line-colouring reduction}\label{sec:reduction}

For $R,C,D,A\subseteq\Zn$ put
\[
\Bl(R,C,D,A)=\{(r,c)\in G: r\in R,\ c\in C,\ r-c\in D,\ r+c\in A\},
\]
and let $\Wh(R,C,D,A)=\Bl(R^{c},C^{c},D^{c},A^{c})$, where ${}^{c}$ denotes complement in $\Zn$ taken family by family.

\begin{lemma}[Line-colouring equivalence]\label{lem:lines}
Let $m\ge 1$. There is a peaceable placement of $m$ black and $m$ white queens on the $15\times 15$ torus if and only if there are sets $R,C,D,A\subseteq\Zn$ with
\[
|\Bl(R,C,D,A)|\ge m
\quad\text{and}\quad
|\Wh(R,C,D,A)|\ge m .
\]
\end{lemma}

\begin{proof}
($\Rightarrow$) Given a peaceable placement, colour every line carrying a black queen black and every line carrying a white queen white. No line receives both colours: a line carrying queens of both colours would exhibit a black--white attack. Colour every remaining line white, and let $R,C,D,A$ be the black line sets in the four families; the white line sets are then exactly the complements $R^{c},C^{c},D^{c},A^{c}$. Every black queen sits on four black lines, so its cell lies in $\Bl$, and the $m$ black queens occupy distinct cells, whence $|\Bl|\ge m$; likewise every white queen sits on four lines that carry no black queen, hence on four white lines, so $|\Wh|\ge m$.

($\Leftarrow$) Given such sets, place black queens on any $m$ cells of $\Bl$ and white queens on any $m$ cells of $\Wh$. The two sets are disjoint, since a cell cannot have all four of its lines both in and out of the black sets. If a black queen attacked a white queen they would share a line, say a row $r$; but then $r\in R$ and $r\in R^{c}$, which is absurd.
\end{proof}

Colouring the unused lines white rather than leaving them unused is harmless in the direction that matters: it can only enlarge the white sets, and $\Wh$ is monotone increasing in the complements. From now on, a \emph{configuration} is a quadruple $(R,C,D,A)$ of subsets of $\Zn$, and we ask whether one exists with $|\Bl|\ge 21$ and $|\Wh|\ge 21$. Write
\[
s_0=|R|,\quad s_1=|C|,\quad s_2=|D|,\quad s_3=|A|,\qquad S=s_0+s_1+s_2+s_3
\]
and call $(s_0,s_1,s_2,s_3)$ the \emph{size profile} of the configuration.

\begin{lemma}[Budget]\label{lem:budget}
If some configuration has $|\Bl|\ge 21$ and $|\Wh|\ge 21$, then some configuration with the same property has $S\le 30$.
\end{lemma}

\begin{proof}
Replacing $(R,C,D,A)$ by $(R^{c},C^{c},D^{c},A^{c})$ exchanges $\Bl$ and $\Wh$, hence preserves the property, and replaces $S$ by $60-S$. One of $S$ and $60-S$ is at most $30$.
\end{proof}

\begin{remark}\label{rem:notequal}
Lemma~\ref{lem:budget} is an inequality and must be kept as one. It is tempting to argue that one may add lines to the black side until $S=30$; this is false, because a line added to $R$ is removed from $R^{c}$ and can destroy $|\Wh|\ge 21$. The witness of Proposition~\ref{prop:lower} illustrates the point: its black side has $S=28$ and its white side $S=32$, and neither is $30$. Accordingly the search below covers the whole range $20\le S\le 30$, not merely the top stratum.
\end{remark}

\section{Admissible size profiles}\label{sec:profiles}

\begin{lemma}[Pair bounds]\label{lem:pairs}
Let $(R,C,D,A)$ be a configuration with size profile $(s_0,s_1,s_2,s_3)$. Then for all $i<j$,
\[
|\Bl|\le s_is_j,\qquad |\Wh|\le (15-s_i)(15-s_j).
\]
Consequently, if $|\Bl|\ge 21$ and $|\Wh|\ge 21$ then
\begin{equation}\label{eq:pair}
s_is_j\ge 21\quad\text{and}\quad (15-s_i)(15-s_j)\ge 21\qquad\text{for all }i<j .
\end{equation}
\end{lemma}

\begin{proof}
Two lines from distinct families meet in exactly one cell. For a row and a column this is immediate; for a row $r$ and a diagonal $d$ the unique common cell is $(r,r-d)$, and similarly for the three remaining mixed pairs involving a row or a column; for a diagonal $d$ and an antidiagonal $a$ it is $(8(d+a),8(a-d))$ by \eqref{eq:inverse}. Every cell of $\Bl$ lies on a line of family $i$ belonging to the black set and on a line of family $j$ belonging to the black set, and distinct cells of $\Bl$ give distinct such pairs of lines; hence $|\Bl|\le s_is_j$. The bound for $\Wh$ is the same statement applied to the complementary sets.
\end{proof}

\begin{lemma}[Triple identity]\label{lem:triples}
Let $X,Y,Z$ be three of the four families, with black sets of sizes $x,y,z$. Let $T$ be the number of cells whose lines in these three families are all black, and $T'$ the number whose lines in these three families are all white. Then
\begin{equation}\label{eq:triple-id}
T+T'=225-15(x+y+z)+xy+xz+yz .
\end{equation}
Consequently, if $|\Bl|\ge 21$ and $|\Wh|\ge 21$ then
\begin{equation}\label{eq:triple}
225-15(x+y+z)+xy+xz+yz\ge 42
\end{equation}
for each of the four triples of families.
\end{lemma}

\begin{proof}
Write $u,v,w\colon G\to\{0,1\}$ for the indicators of ``the cell's line in family $X$ (resp.\ $Y$, $Z$) is black''. Then
\[
uvw+(1-u)(1-v)(1-w)=1-u-v-w+uv+uw+vw .
\]
Summing over the $225$ cells: $\sum 1=225$; $\sum u=15x$, since each of the $x$ black lines of family $X$ contains $15$ cells, and likewise $\sum v=15y$, $\sum w=15z$; and $\sum uv=xy$, since each pair consisting of a black line of family $X$ and a black line of family $Y$ contributes exactly one cell, by the incidence fact used in Lemma~\ref{lem:pairs}; likewise $\sum uw=xz$ and $\sum vw=yz$. This is \eqref{eq:triple-id}. Every cell of $\Bl$ is counted by $T$ and every cell of $\Wh$ by $T'$, and the two sets are disjoint, so $|\Bl|+|\Wh|\le T+T'$; if both are at least $21$ we get \eqref{eq:triple}.
\end{proof}

Note that \eqref{eq:triple-id} and hence \eqref{eq:triple} are symmetric under $(x,y,z)\mapsto(15-x,15-y,15-z)$, as they must be.

\subsection*{Symmetry}

Let $\Gamma$ be the group of affine bijections of $G$ that permute the four line families. It consists of the $225$ translations composed with $64$ linear maps, so $|\Gamma|=14{,}400$; this was verified by exhaustive enumeration of the invertible $2\times2$ matrices over $\Zn$. Every $\gamma\in\Gamma$ carries a configuration $(R,C,D,A)$ to a configuration with the same $|\Bl|$ and $|\Wh|$, permuting the families according to a homomorphism $\Gamma\to\mathrm{Sym}\{R,C,D,A\}$. The image is the dihedral group of order $8$ preserving the partition $\{\{R,C\},\{D,A\}\}$, generated by
\begin{align*}
\text{transpose } \tau(r,c)&=(c,r): && R\leftrightarrow C,\ D\mapsto -D,\ A\mapsto A;\\
\text{reflection } \nu(r,c)&=(r,-c): && R\mapsto R,\ C\mapsto -C,\ D\leftrightarrow A;\\
\text{halving } \psi(r,c)&=(r+c,\ r-c): && \{R,C\}\leftrightarrow\{D,A\} .
\end{align*}
For $\psi$ one checks directly that a level set of the new row coordinate is an old antidiagonal, of the new column coordinate an old diagonal, of the new diagonal an old column and of the new antidiagonal an old row; $\psi$ is bijective because $\det\binom{1\ \ 1}{1\ -1}=-2$ is a unit modulo $15$. The image is not all of $\mathrm{Sym}_4$: the four line directions $\infty,0,1,-1$ of $\mathbb{P}^1(\Zn)$ have harmonic cross-ratio, and the family action is the corresponding dihedral group of order $8$. The induced action on size profiles is the permutation action of this $D_4$ on ordered quadruples.

\begin{proposition}\label{prop:247}
Let $(R,C,D,A)$ be a configuration with $|\Bl|\ge 21$, $|\Wh|\ge 21$ and $S\le 30$. Then its size profile is equivalent, under the $D_4$ action together with complementation $s_i\mapsto 15-s_i$ in the case $S=30$, to one of exactly $247$ canonical profiles, distributed by $S$ as in Table~\ref{tab:strata}.
\end{proposition}

\begin{table}[h]
\centering
\begin{tabular}{@{}l*{11}{r}@{}}
\toprule
$S$ & 20 & 21 & 22 & 23 & 24 & 25 & 26 & 27 & 28 & 29 & 30 \\
canonical profiles & 1 & 1 & 4 & 6 & 13 & 18 & 29 & 37 & 47 & 52 & 39 \\
\bottomrule
\end{tabular}
\medskip
\caption{The $247$ canonical size profiles by budget $S$.}\label{tab:strata}
\end{table}

\begin{proof}
Filter all $16^4$ quadruples $(s_0,s_1,s_2,s_3)\in\{0,\dots,15\}^4$ by $S\le 30$ and by conditions \eqref{eq:pair} and \eqref{eq:triple}, which are necessary by Lemmas~\ref{lem:pairs} and~\ref{lem:triples}, and canonicalize the survivors: replace each by the lexicographically least element of its $D_4$ orbit, adjoining the $D_4$ orbit of the complementary profile $(15-s_i)_i$ when $S=30$. Complementation preserves the property ``$|\Bl|\ge21$ and $|\Wh|\ge21$'' but sends $S$ to $60-S$, so it may be used as an identification only within the stratum $S=30$, which it preserves. The enumeration is a finite computation over $65{,}536$ quadruples; it was carried out independently three times, in two languages, always giving the same $247$ profiles and the distribution of Table~\ref{tab:strata}.
\end{proof}

Inspection of the list shows that every surviving profile satisfies $3\le s_i\le 12$ for all $i$; the unique profiles at the two smallest budgets are $(5,5,5,5)$ and $(5,5,5,6)$, while the top stratum $S=30$ contains both the lopsided $(3,8,9,10)$ and the balanced $(7,7,8,8)$.

\section{Deciding one profile exactly}\label{sec:completion}

Fix a canonical profile. By the $D_4$ symmetry we may choose either $(R,C)$ or $(D,A)$ as the \emph{fixed pair}; write $(r,c)$ for its two sizes, with $r\le c$ after a transpose if necessary, and $(d,a)$ for the sizes of the two remaining families, which we call the \emph{completion} families. The search proceeds in three nested stages: orbit representatives for the fixed pair, all subsets for one completion family, and an exact decision for the last.

\subsection*{Stage 1: representatives of the fixed pair}

Fix the pair $(R,C)$ up to the subgroup of $\Gamma$ that preserves the ordered profile. This subgroup is generated by:
\begin{itemize}
\item independent translations $(R,C)\mapsto(R+\alpha,\ C+\beta)$, which change $D$ and $A$ only by translations ($225$ elements);
\item common unit scaling $(R,C)\mapsto(uR,\ uC)$ for $u\in(\Zn)^{\times}$, which scales $D$ and $A$ by the same $u$ ($8$ elements);
\item the transpose $\tau$, admissible exactly when $r=c$, since $\tau$ exchanges the two fixed families while sending $D\mapsto -D$ and fixing $A$, hence preserving the completion sizes;
\item the relative sign $\nu$, i.e.\ $(R,C)\mapsto(R,-C)$, admissible \emph{only} when $d=a$.
\end{itemize}
The last restriction is the one soundness point that is easy to get wrong. Since $\nu$ exchanges the diagonal and antidiagonal families, an ordered profile with $d\ne a$ is not mapped to itself, and quotienting by $\nu$ would then identify configurations belonging to different ordered profiles --- legitimate only if both orientations are enumerated. We use $\nu$ only when $d=a$, where it is a genuine automorphism of the ordered profile. An earlier cost model in this project applied $\nu$ unconditionally; correcting it raised the workload of the independent enumerator from $4{,}876{,}723{,}839$ to $6{,}241{,}793{,}402$ completions.

The resulting group has order $225\cdot 8=1800$, doubled for each of the two admissible extra generators, hence $1800$, $3600$ or $7200$. Representatives are produced by first reducing each set to its cyclic-translation necklace and then quotienting by the residual unit, transpose and sign actions. Only $43$ distinct triples $(r,c,\text{group})$ occur across the $247$ profiles, and every one of the resulting orbit counts was verified independently by Burnside's lemma (Section~\ref{sec:verification}).

Choosing the fixed pair is a free optimization: for each profile both choices are costed as (number of pair representatives) $\times$ (number of subsets of the smaller completion family) and the cheaper one is taken. Likewise, within the completion families, the family with fewer subsets is the one enumerated in Stage~2, which is legitimate because $\nu$ exchanges the two completion families globally, so the ordered profiles $(r,c,d,a)$ and $(r,c,a,d)$ are simultaneously feasible or infeasible.

\subsection*{Stage 2: the completion as a subset problem}

Fix a representative $(R,C)$. Index cells by their diagonal and antidiagonal coordinates via \eqref{eq:inverse}, writing $x$ for the diagonal and $y$ for the antidiagonal coordinate to keep them apart from the sizes $d,a$, and set
\[
P_{x,y}=\mathbf 1_R\bigl(8(x+y)\bigr)\,\mathbf 1_C\bigl(8(y-x)\bigr),
\qquad
Q_{x,y}=\mathbf 1_{R^{c}}\bigl(8(x+y)\bigr)\,\mathbf 1_{C^{c}}\bigl(8(y-x)\bigr)
\qquad (x,y\in\Zn).
\]
For each antidiagonal set $A\subseteq\Zn$ of the prescribed size $a$, put
\[
p_x=\sum_{y\in A}P_{x,y},
\qquad
q_x=\sum_{y\notin A}Q_{x,y},
\qquad
Q_0=\sum_{x\in\Zn}q_x .
\]
Then for any diagonal set $D$,
\[
|\Bl(R,C,D,A)|=\sum_{x\in D}p_x,
\qquad
|\Wh(R,C,D,A)|=\sum_{x\notin D}q_x=Q_0-\sum_{x\in D}q_x .
\]
Hence a diagonal set $D$ of the prescribed size $d$ completes the configuration if and only if
\begin{equation}\label{eq:subset}
\sum_{x\in D}p_x\ \ge\ 21
\qquad\text{and}\qquad
\sum_{x\in D}q_x\ \le\ Q_0-21 .
\end{equation}
This is an exact $15$-item, fixed-cardinality, two-constraint subset problem, with nonnegative integer weights $p_x,q_x\le 15$. Before it is posed, the necessary conditions $\sum_x p_x\ge 21$ and $Q_0\ge 21$ are tested and cheaply discard most candidates $A$.

\subsection*{Stage 3: exact decision of \eqref{eq:subset}}

Problem \eqref{eq:subset} is decided by a depth-first search over the $15$ items, each included or excluded, with the following bounds precomputed by suffix dynamic programming: for every position $x$ and every remaining cardinality $k$, the largest attainable sum of $k$ of the weights $p_x,\dots,p_{14}$, and the smallest attainable sum of $k$ of the weights $q_x,\dots,q_{14}$.

\begin{lemma}\label{lem:exhaustive}
The search of Stage~3 answers \eqref{eq:subset} correctly.
\end{lemma}

\begin{proof}
A node is abandoned only when one of the following holds: the number of remaining items is smaller than the cardinality still required; the accumulated $q$-sum already exceeds $Q_0-21$; the accumulated $p$-sum plus the maximal attainable remaining $p$-sum with the required cardinality is below $21$; or the accumulated $q$-sum plus the minimal attainable remaining $q$-sum with the required cardinality exceeds $Q_0-21$. Each is a necessary condition for a feasible completion within the subtree, so no feasible $D$ is discarded. Conversely the search reports success only at a node where the required cardinality is reached and the $p$-constraint is met, or at a node where the $p$-constraint is already met and the minimal remaining $q$-sum with the required cardinality keeps the $q$-constraint satisfiable; in the latter case an explicit feasible $D$ exists, since $p$-sums only increase as further items are added. In both cases a feasible $D$ exists. As the search is otherwise a complete binary recursion over the $15$ items, it is exhaustive.
\end{proof}

Item order affects only speed, not correctness; the implementation sorts items by a ratio heuristic. Consequently the three stages together decide, for a given profile, whether \emph{any} configuration with that ordered size profile has $|\Bl|\ge 21$ and $|\Wh|\ge 21$, and the answer transfers to the whole $D_4$ orbit and, in the stratum $S=30$, to the complementary profile as well.

\section{The search and its outcome}\label{sec:result}

\begin{proposition}\label{prop:unsat}
For every one of the $247$ canonical profiles of Proposition~\ref{prop:247}, no configuration with that profile has $|\Bl|\ge 21$ and $|\Wh|\ge 21$.
\end{proposition}

The proof is the exhaustive execution of Section~\ref{sec:completion} on all $247$ profiles. The run statistics are given in Table~\ref{tab:stats}; a per-profile certificate recording, for each profile, the fixed pair, the orbit count, a hash of the representative list, and the case counts, is included among the ancillary files.

\begin{table}[h]
\centering
\begin{tabular}{@{}lr@{}}
\toprule
canonical profiles                 & $247$ \\
distinct fixed-pair orbit types    & $43$ \\
antidiagonal subsets examined      & $6{,}074{,}753{,}568$ \\
survivors of the necessary-condition prefilter & $2{,}329{,}772{,}941$ \\
calls to the exact subset decision & $26{,}948{,}453$ \\
branch nodes in those calls        & $246{,}476{,}279$ \\
profiles admitting a $21+21$ configuration & $0$ \\
\bottomrule
\end{tabular}
\medskip
\caption{Statistics of the exhaustive run. The number of antidiagonal subsets examined equals $\sum_{\text{profiles}} (\text{orbit count})\cdot\binom{15}{a}$, which is checked row by row by the audit script.}\label{tab:stats}
\end{table}

\begin{proof}[Proof of Theorem~\ref{thm:main}]
The lower bound is Proposition~\ref{prop:lower}. For the upper bound, suppose $21$ black and $21$ white queens are placed peaceably. By Lemma~\ref{lem:lines} there is a configuration with $|\Bl|\ge 21$ and $|\Wh|\ge 21$, by Lemma~\ref{lem:budget} one may take $S\le 30$, and by Proposition~\ref{prop:247} its size profile lies in the $247$-element canonical list. This contradicts Proposition~\ref{prop:unsat}. Hence $t(15)\le 20$, and $t(15)=20$.
\end{proof}

\section{Verification}\label{sec:verification}

An exhaustive computation is only as good as the evidence that it is exhaustive and correctly implemented. We list the checks performed. The two enumerators mentioned below --- referred to as the \emph{proof solver} and the \emph{independent enumerator} --- were written separately, from the mathematical description of Section~\ref{sec:completion}, in different C++ dialects, by different systems (see Section~\ref{sec:ai}), and share no code.

\emph{Two independent enumerators.} The proof solver (\texttt{peace15\_solver.cpp}) and the independent enumerator (\texttt{profile\_enum.cpp}) agree that all $247$ profiles are infeasible. They differ in three substantive respects. First, in the symmetry quotient: the two select different fixed pairs on $19$ of the $247$ profiles, and consequently enumerate different case sets --- $6{,}074{,}753{,}568$ antidiagonal subsets for the solver against $6{,}241{,}793{,}402$ completions for the enumerator. Second, in the decision procedure for \eqref{eq:subset}: the solver uses the depth-first search with suffix dynamic-programming bounds of Stage~3, whereas the enumerator uses an exact $7+8$ meet-in-the-middle table, retaining for each cardinality and each $p$-sum threshold the minimum achievable $q$-sum over right-half subsets --- an algorithm with no search and no pruning at all. Third, in scheduling: the solver is OpenMP-parallel, the enumerator single-threaded and resumable, writing one JSON record per profile.

\emph{Agreement on the profile list.} The $247$-profile worklist was generated independently by us in Python, from Lemmas~\ref{lem:pairs} and~\ref{lem:triples}; the proof solver regenerates the set from scratch in C++ at startup and aborts unless it coincides with the supplied file, which it does byte for byte; the audit script regenerates it a third time in Python. At $S=30$ the count was also reproduced by an independent hand-checkable enumeration as $370$ ordered profiles falling into $39$ canonical classes.

\emph{Repeated runs on different toolchains.} The proof solver was run three times: twice as recorded by its author (five threads, $60.4$~s; three threads, $74.0$~s) and once by us from a fresh build with a different compiler and no OpenMP (single-threaded). All $247$ rows of the certificate agree in every column except the per-profile timing.

\emph{Brute-force gates on the completion kernel.} At startup the proof solver compares its Stage~3 decision against direct enumeration of all $\binom{15}{k}$ subsets on $20{,}000$ deterministic pseudorandom instances with randomized weights and cardinalities, and aborts on any disagreement; the independent enumerator runs the same style of gate on $100{,}000$ instances. The proof solver additionally checks three known pair-orbit counts, $11{,}793$, $6{,}892$ and $13{,}654$, before starting.

\emph{Burnside audit.} An audit script (\texttt{peace15\_audit.py}) recomputes all $43$ fixed-pair orbit counts from Burnside's lemma --- counting, for each group element, the subsets of each cardinality invariant under the induced affine permutation of $\Zn$, via the cycle structure --- and compares them with the counts recorded in the certificate. It also re-derives the profile list, verifies that each certificate row satisfies $\text{cases}=\text{orbits}\times\binom{15}{a}$, verifies the column total $6{,}074{,}753{,}568$, and re-counts the witness of Proposition~\ref{prop:lower} as $20$ black-homogeneous and $22$ white-homogeneous cells.

\emph{Memory safety.} An AddressSanitizer and UndefinedBehaviorSanitizer build of the proof solver passes all startup self-tests.

\emph{The witness.} The $20+20$ placement of Proposition~\ref{prop:lower} was checked by four independently written scripts, including a from-scratch check that verifies all $400$ black--white pairs against all four attack relations directly, rather than through the line-set formulation.

\emph{Source audit.} The proof solver was read line by line against the mathematics of Section~\ref{sec:completion}, with particular attention to the three places where an unsound shortcut would be invisible in the output: the admissibility conditions for the transpose and relative-sign quotients, the necessary-condition character of every pruning rule, and the reconstruction path that would turn a reported feasible instance into an explicit quadruple verified by direct enumeration of the $225$ cells.

\emph{Independent partial evidence.} A SAT-based cube-and-conquer campaign on a $3{,}244$-variable propositional encoding of the $21+21$ question, later superseded and never finished, produced roughly $28{,}900$ DRAT-certified cube refutations checked by \texttt{drat-trim} in default backward mode, with no satisfiable cube anywhere. It proves nothing on its own but is consistent with Theorem~\ref{thm:main}. A per-profile certified-CNF pilot refuted $(5,5,5,5)$ with a checked $9.2$~MB DRAT proof but timed out on $(6,8,7,7)$ under two solvers; this is why the exhaustive enumeration, and not a proof-logging pipeline, carries the upper bound here.

\emph{What is not claimed.} The certificate files are audit logs, not standalone proof objects in the sense of DRAT or LRAT: they record what was searched, not a machine-checkable refutation. The refutation is the short program together with the reductions of Sections~\ref{sec:reduction}--\ref{sec:completion}; re-running the program is what re-establishes it, and agreement between two implementations is a strong error filter, not a formal guarantee.

\section{Reproduction and ancillary files}\label{sec:repro}

The ancillary files accompanying this note are:
\begin{center}
\begin{tabular}{@{}ll@{}}
\toprule
\texttt{peace15\_solver.cpp} & the proof solver (C++17, OpenMP)\\
\texttt{peace15\_audit.py} & profile, Burnside and witness audit (Python 3)\\
\texttt{profile\_enum.cpp} & the independent enumerator (C++20)\\
\texttt{WORKLIST-2026-07-25.tsv} & the $247$ canonical profiles\\
\texttt{peace15\_certificate.tsv} & per-profile certificate, five-thread run\\
\texttt{peace15\_certificate\_run2.tsv} & per-profile certificate, three-thread run\\
\texttt{README.txt} & build and run instructions, expected output\\
\bottomrule
\end{tabular}
\end{center}

To rebuild the upper bound:
\begin{verbatim}
g++ -O3 -march=native -fopenmp -DNDEBUG -std=c++17 \
    peace15_solver.cpp -o peace15_solver
OMP_NUM_THREADS=5 ./peace15_solver \
    --worklist WORKLIST-2026-07-25.tsv \
    --output peace15_certificate.tsv
python3 peace15_audit.py
\end{verbatim}
The solver prints \texttt{SELFTEST\_OK} after its startup gates and \texttt{ALL\_UNSAT} on completion; the audit script prints \texttt{AUDIT\_OK}. The recorded five-thread run took about $60$ seconds; runtime is not part of the proof. The independent enumerator is built and run by
\begin{verbatim}
clang++ -O3 -std=c++20 -DNDEBUG profile_enum.cpp -o profile_enum
./profile_enum --self-test --random 100000
./profile_enum --batch --worklist WORKLIST-2026-07-25.tsv \
    --results results.jsonl --log logs
\end{verbatim}
and takes about $200$ seconds single-threaded.

SHA-256 digests of the two files that carry the upper bound, as recorded for the first run:
\begin{center}\small
\begin{tabular}{@{}l@{}}
\toprule
\texttt{peace15\_solver.cpp}\\
\quad\texttt{1ae5b340d6feb4591dae4f7b57a283a373ce60fd2bc9767fc043f25d402bebe9}\\
\addlinespace
\texttt{peace15\_certificate.tsv}\\
\quad\texttt{f9d2c858cc6e3a144a10da509a744114080de801df6d216b8cd2783cb2f410b0}\\
\bottomrule
\end{tabular}
\end{center}

\section{Use of AI tools}\label{sec:ai}

This note was produced with substantial use of AI systems, which we disclose in the interest of transparency. The reduction programme --- the line-colouring equivalence, the budget inequality, the pair and triple bounds, the symmetry analysis and the $247$-profile worklist --- was developed inside an AI-assisted research repository in which Anthropic Claude agents ran the computations and maintained the standing documents, and in which a large-language-model collaborator acting as a mathematician was consulted in writing; the exact completion algorithm of Section~\ref{sec:completion} originates in that consultation. One error in that reply (a claim that the budget could be normalized to $S=30$) was found and corrected here; see Remark~\ref{rem:notequal}. The final proof document, the proof solver \texttt{peace15\_solver.cpp} and the audit script \texttt{peace15\_audit.py} were produced by OpenAI GPT-5.6~Pro on July~25, 2026.

The verification described in Section~\ref{sec:verification} was carried out by Anthropic Claude agents: the second exact enumerator \texttt{profile\_enum.cpp} and its brute-force gates, the line-by-line audit of the proof solver, the local rebuild and rerun on a different toolchain, the sanitizer builds, the re-derivation and numerical checking of every identity in Sections~\ref{sec:reduction}--\ref{sec:profiles}, and the independent checks of the witness. The soundness correction to the fixed-pair quotient described in Section~\ref{sec:completion} was found during that verification.

Agreement between two AI-written implementations is a useful and, here, deliberately engineered error filter, but it is not independent verification in the scholarly sense, and it does not replace the proofs given in the text. Readers who wish to check Theorem~\ref{thm:main} should re-run the ancillary programs or, better, write a third implementation from Sections~\ref{sec:reduction}--\ref{sec:completion}. The author reviewed the manuscript and takes full responsibility for its content.

\section*{Acknowledgements}

The problem was brought within reach by the work of Clinch, Drescher, Huynh and Saffidine, whose odd-torus bound bracketed $t(15)$ from above and whose construction supplies the matching lower bound, and by the OEIS contributors who have kept A279405 current --- most recently Jaideep Padhi, whose values for $n=13$ and $n=14$ in June 2026 marked $n=15$ as the live frontier.

\begin{thebibliography}{99}

\bibitem{Ainley77}
S.~Ainley,
\emph{Mathematical Puzzles},
G.~Bell \& Sons, London, 1977.

\bibitem{CDHS24}
S.~Clinch, M.~Drescher, T.~Huynh, and A.~Saffidine,
\emph{Constructions, bounds, and algorithms for peaceable queens},
arXiv:2406.06974 (2024).

\bibitem{OEIS250000}
N.~J.~A. Sloane and others,
\emph{The On-Line Encyclopedia of Integer Sequences}, sequence A250000:
\emph{Peaceable coexisting armies of queens},
\url{https://oeis.org/A250000}, accessed July 25, 2026.

\bibitem{OEIS279405}
N.~J.~A. Sloane and others,
\emph{The On-Line Encyclopedia of Integer Sequences}, sequence A279405:
\emph{Peaceable coexisting armies of queens on a toroidal board},
\url{https://oeis.org/A279405}, accessed July 25, 2026.

\end{thebibliography}

\end{document}
